\documentclass[11pt, a4paper]{article}

\usepackage[utf8]{inputenc}
\usepackage[T1]{fontenc}
\usepackage[english]{babel}
\usepackage{amsmath, amssymb}
\usepackage{geometry}
\usepackage{booktabs}
\usepackage{graphicx}
\usepackage{xcolor}
\usepackage[colorlinks=true, linkcolor=blue!60!black, urlcolor=blue!60!black]{hyperref}
\usepackage{microtype}
\usepackage{enumitem}
\usepackage{parskip}
\usepackage{titlesec}
\usepackage{listings}

\geometry{margin=2.4cm}

\titleformat{\section}{\large\bfseries}{\thesection.}{0.6em}{}
\titleformat{\subsection}{\normalsize\bfseries}{\thesubsection.}{0.5em}{}

\lstset{
    basicstyle=\ttfamily\small,
    keywordstyle=\color{blue!70!black}\bfseries,
    commentstyle=\color{gray},
    frame=single,
    breaklines=true,
    columns=fullflexible,
    language=Python,
}

\title{\textbf{Bayer superpixel weight derivation}\\
       \large Thorlabs LP126CU camera under Olympus U-LH100IR illumination}
\author{InterferoLab}
\date{May 2026}

\begin{document}

\maketitle
\tableofcontents
\vspace{1em}

% ══════════════════════════════════════════════════════════
\section{Motivation}

The Thorlabs LP126CU is a color CMOS sensor with a \textbf{Bayer color
filter array} (CFA) in RGGB pattern\@. Each physical pixel measures only
one of the three colors, according to its position in the $2\times 2$
grid:

\begin{center}
\begin{tabular}{|c|c|}
\hline
R & G\textsubscript{1} \\
\hline
G\textsubscript{2} & B \\
\hline
\end{tabular}
\end{center}

Coherence scanning interferometry (CSI) works with a \emph{monochrome}
intensity signal: the gray value of each position $(y,x)$ as a function
of the axial position $z$. To exploit the signal of all colors without
resorting to demosaicing (which introduces spurious harmonics, see the
\emph{Camera characterization report}), each $2\times 2$ block is
collapsed into a single \textbf{superpixel}:
\begin{equation}
  I_{\mathrm{sp}} = w_R\,R + w_{G_1}\,G_1 + w_{G_2}\,G_2 + w_B\,B.
  \label{eq:superpixel}
\end{equation}
The optimal choice of the weights $w_c$ maximizes the signal-to-noise
ratio (SNR) of the superpixel.

% ══════════════════════════════════════════════════════════
\section{Optimal weighting criterion}

The CSI interference signal is proportional to the received light
intensity. For each channel $c \in \{R, G_1, G_2, B\}$, the captured
signal is:
\begin{equation}
  S_c = \int_{\lambda} S_{\mathrm{lamp}}(\lambda)\;
                       T_{\mathrm{IR}}(\lambda)\;
                       \mathrm{QE}_c(\lambda)\;d\lambda
  \label{eq:signal}
\end{equation}
where:
\begin{itemize}[noitemsep]
  \item $S_{\mathrm{lamp}}(\lambda)$ — emission spectrum of the lamp
        (spectral power density, a.u.).
  \item $T_{\mathrm{IR}}(\lambda)$ — transmission of the IR filter built
        into the camera.
  \item $\mathrm{QE}_c(\lambda)$ — quantum efficiency of channel $c$
        (taken from the Thorlabs datasheet, without the IR filter, which
        is applied separately).
\end{itemize}

The optimal weight in the maximum-SNR sense is simply proportional to
the received signal:
\begin{equation}
  w_c \propto S_c, \qquad \sum_c w_c = 1.
  \label{eq:weights}
\end{equation}
The RGGB block contains \emph{two} green pixels for every red and blue
one, so the normalization denominator is $S_R + 2S_G + S_B$, not
$S_R + S_G + S_B$:
\begin{equation}
  w_R = \frac{S_R}{S_R + 2S_G + S_B}, \quad
  w_{G_1} = w_{G_2} = \frac{S_G}{S_R + 2S_G + S_B}, \quad
  w_B = \frac{S_B}{S_R + 2S_G + S_B}.
  \label{eq:normweights}
\end{equation}

% ══════════════════════════════════════════════════════════
\section{Spectral model of the lamp}

The lamp in use is an \textbf{Olympus U-LH100IR} (100\,W, 12\,V), a
tungsten halogen lamp. Halogen lamps operate at higher filament
temperatures (${\sim}3200$\,K) than conventional incandescent ones
(${\sim}2800$\,K) thanks to the halogen cycle that regenerates the
filament.

The emission spectrum is approximated by \textbf{Planck's law} for a
black body at $T = 3200$\,K:
\begin{equation}
  B(\lambda, T) = \frac{2hc^2}{\lambda^5}
                  \cdot \frac{1}{e^{hc/(\lambda k_B T)} - 1}
\end{equation}
Only the spectral shape matters, not the absolute magnitude.

\textbf{Note:} If the measured spectral curve of the U-LH100IR is
available, $B(\lambda, T)$ should be replaced by that curve for better
accuracy.

% ══════════════════════════════════════════════════════════
\section{Input data}

The script \texttt{compute\_bayer\_weights.py} reads four CSV files
extracted from the plots in the Thorlabs datasheet with the
WebPlotDigitizer tool:

\begin{center}
\begin{tabular}{ll}
\toprule
\textbf{File} & \textbf{Contents} \\
\midrule
\texttt{ir\_filter\_transmission.csv} & $T_{\mathrm{IR}}(\lambda)$ in \% \\
\texttt{qe\_red.csv}   & $\mathrm{QE}_R(\lambda)$ in \% \\
\texttt{qe\_green.csv} & $\mathrm{QE}_G(\lambda)$ in \% \\
\texttt{qe\_blue.csv}  & $\mathrm{QE}_B(\lambda)$ in \% \\
\bottomrule
\end{tabular}
\end{center}

All data are interpolated onto a common 400--1000\,nm grid with a 1\,nm
step. Values outside the measured range are treated as zero (null
sensor response), which is the correct conservative assumption.

% ══════════════════════════════════════════════════════════
\section{Computation procedure}

\begin{enumerate}
  \item \textbf{Load and interpolate} the four curves onto the grid
        $\lambda \in [400, 1000]$\,nm.
  \item \textbf{Compute the lamp spectrum} $B(\lambda, 3200)$ from
        Planck's law.
  \item \textbf{Compute the integrand} per channel:
        \begin{equation}
          f_c(\lambda) = B(\lambda, 3200)\cdot T_{\mathrm{IR}}(\lambda)
                        \cdot \mathrm{QE}_c(\lambda)
        \end{equation}
  \item \textbf{Integrate numerically} with the trapezoidal rule:
        $S_c = \int f_c(\lambda)\,d\lambda$.
  \item \textbf{Normalize} according to equation~\eqref{eq:normweights}:
        denominator $= S_R + 2S_G + S_B$.
\end{enumerate}

The script also generates a verification figure with three panels:
(a) the interpolated input curves, (b) the normalized lamp spectrum,
and (c) the integrands $f_c(\lambda)$ normalized by their own maximum.

% ══════════════════════════════════════════════════════════
\section{Results}

With $T = 3200$\,K (Olympus U-LH100IR) and the LP126CU curves:

\begin{center}
\begin{tabular}{lrr}
\toprule
\textbf{Channel} & \textbf{Weight} & \textbf{Integrated signal (a.u.)} \\
\midrule
$w_R$             & 0.257090 & $S_R$ \\
$w_{G_1}$         & 0.311272 & $S_G$ \\
$w_{G_2}$         & 0.311272 & $S_G$ \\
$w_B$             & 0.120366 & $S_B$ \\
\midrule
Sum               & 1.000000 & \\
\bottomrule
\end{tabular}
\end{center}

The green channel receives the largest individual weight because its QE
is high across the visible range and the halogen lamp emits strongly in
that region. The blue channel is the weakest: the QE is low ($<$ 20\%)
in the blue, and the lamp spectrum decays exponentially towards short
wavelengths.

\subsection{Comparison with 2800\,K}

\begin{center}
\begin{tabular}{lcccc}
\toprule
$T$ (K) & $w_R$ & $w_{G_1}=w_{G_2}$ & $w_B$ & Change in $w_B$ \\
\midrule
2800 & 0.2829 & 0.3073 & 0.1026 & — \\
3200 & 0.2571 & 0.3113 & 0.1204 & $+17\%$ \\
\bottomrule
\end{tabular}
\end{center}

The correct temperature (3200\,K vs.\ 2800\,K) increases the blue
channel's weight by 17\% and reduces the red one's by 9\%. Using
2800\,K would over-weight red (which carries less real signal from the
halogen lamp) and under-weight blue.

% ══════════════════════════════════════════════════════════
\section{Use of the weights in the system}

Since the July 2026 reorganization, the \textbf{single source} of the
weights on the Python side is the module
\texttt{utils/camera\_constants.py}. It is consumed by:

\begin{enumerate}
  \item \textbf{Acquisition} (\texttt{backend/acquisition/acquisition\_controller.py}):
        the function \texttt{\_bayer\_to\_superpixel()}, which combines
        the four sub-pixels of the raw Bayer block in real time during
        acquisition.

  \item \textbf{Offline processing} (\texttt{scripts/process\_superpixel.py}):
        conversion of already-acquired datasets to the superpixel format.

  \item \textbf{GUI}: the default values of the weight spinboxes in the
        processing panel.
\end{enumerate}

The \textbf{C++ analysis} keeps a copy of these values in
\texttt{backend/analysis/include/config.hpp} (constants
\texttt{GRAY\_WEIGHT\_R/G/B}), which must be kept \emph{synchronized by
hand} with \texttt{utils/camera\_constants.py} on recalibration (a
regression test detects divergence). In that context the images are
already demosaiced and the G channel represents the mean of $G_1$ and
$G_2$, so its weight is doubled:
\begin{equation}
  W_G^{\mathrm{RGB}} = w_{G_1} + w_{G_2} = 0.622544.
\end{equation}

% ══════════════════════════════════════════════════════════
\section{Limitations and possible improvements}

\begin{itemize}
  \item \textbf{Curves extracted from plots}: the CSVs come from
        WebPlotDigitizer applied to the datasheet plots, not from
        tabulated data. The digitization error is estimated at
        $\pm 1$--$2\%$ per channel.

  \item \textbf{Lamp spectrum}: a black body at 3200\,K is used as an
        approximation. The real color temperature of the U-LH100IR can
        vary slightly with the supply voltage. A measured spectral
        curve would improve accuracy.

  \item \textbf{Optical-path transmission}: the model does not include
        the spectral transmission of the interferometric objective, the
        beam splitter, or the condenser of the Olympus microscope.
        These elements can shift the optimal weights non-negligibly
        (the experimental coherence length suggests a spectral
        narrowing of a factor ${\sim}2$ with respect to the
        camera-only model).
\end{itemize}

\end{document}
